\documentclass[10pt,a4paper]{amsart}
\usepackage{amsmath,amssymb,amsthm,graphicx}
% Retained preamble for the Stacks Project chapter {CH} (shared project preamble).
% The chapter file pinned-original-source.tex only carries \input{{preamble}} on its first line,
% so no preamble text lives inside the chapter itself.  The Stacks Project uses one single shared
% preamble for every chapter: the snapshot stages it three times (sources/stacks-chow/preamble.tex,
% sources/stacks-etale-cohomology/preamble.tex and focused-seeds/AGP/sources/stacks-snapshot-20261011/
% preamble.tex) and those copies are identical in size.  The block below is copied mechanically,
% line for line, from sources/stacks-chow/preamble.tex lines 14-18 and 154-217: the Xy-pic setup and
% the native theorem declarations together with the mathematical macros.  Nothing was added, removed
% or redefined; the wrapper class and the amsmath/amssymb/amsthm/graphicx base packages are supplied
% by the assembler.
% For commutative diagrams we use Xy-pic
% Chapter: stacks-spaces-more-morphisms.tex -- More on Morphisms of Spaces.
\usepackage[all]{xy}

% We use 2cell for 2-commutative diagrams.
\xyoption{2cell}

% Theorem environments.
%
\theoremstyle{plain}
\newtheorem{theorem}[subsection]{Theorem}
\newtheorem{proposition}[subsection]{Proposition}
\newtheorem{lemma}[subsection]{Lemma}

\theoremstyle{definition}
\newtheorem{definition}[subsection]{Definition}
\newtheorem{example}[subsection]{Example}
\newtheorem{exercise}[subsection]{Exercise}
\newtheorem{situation}[subsection]{Situation}

\theoremstyle{remark}
\newtheorem{remark}[subsection]{Remark}
\newtheorem{remarks}[subsection]{Remarks}

\numberwithin{equation}{subsection}

% Macros
%
\def\lim{\mathop{\mathrm{lim}}\nolimits}
\def\colim{\mathop{\mathrm{colim}}\nolimits}
\def\Spec{\mathop{\mathrm{Spec}}}
\def\Hom{\mathop{\mathrm{Hom}}\nolimits}
\def\Ext{\mathop{\mathrm{Ext}}\nolimits}
\def\SheafHom{\mathop{\mathcal{H}\!\mathit{om}}\nolimits}
\def\SheafExt{\mathop{\mathcal{E}\!\mathit{xt}}\nolimits}
\def\Sch{\mathit{Sch}}
\def\Mor{\mathop{\mathrm{Mor}}\nolimits}
\def\Ob{\mathop{\mathrm{Ob}}\nolimits}
\def\Sh{\mathop{\mathit{Sh}}\nolimits}
\def\NL{\mathop{N\!L}\nolimits}
\def\CH{\mathop{\mathrm{CH}}\nolimits}
\def\proetale{{pro\text{-}\acute{e}tale}}
\def\etale{{\acute{e}tale}}
\def\QCoh{\mathit{QCoh}}
\def\Ker{\mathop{\mathrm{Ker}}}
\def\Im{\mathop{\mathrm{Im}}}
\def\Coker{\mathop{\mathrm{Coker}}}
\def\Coim{\mathop{\mathrm{Coim}}}

% Boxtimes
%
\DeclareMathSymbol{\boxtimes}{\mathbin}{AMSa}{"02}

%
% Macros for moduli stacks/spaces
%
\def\QCohstack{\mathcal{QC}\!\mathit{oh}}
\def\Cohstack{\mathcal{C}\!\mathit{oh}}
\def\Spacesstack{\mathcal{S}\!\mathit{paces}}
\def\Quotfunctor{\mathrm{Quot}}
\def\Hilbfunctor{\mathrm{Hilb}}
\def\Curvesstack{\mathcal{C}\!\mathit{urves}}
\def\Polarizedstack{\mathcal{P}\!\mathit{olarized}}
\def\Complexesstack{\mathcal{C}\!\mathit{omplexes}}
% \Pic is the operator that assigns to X its picard group, usage \Pic(X)
% \Picardstack_{X/B} denotes the Picard stack of X over B
% \Picardfunctor_{X/B} denotes the Picard functor of X over B
\def\Pic{\mathop{\mathrm{Pic}}\nolimits}
\def\Picardstack{\mathcal{P}\!\mathit{ic}}
\def\Picardfunctor{\mathrm{Pic}}
\def\Deformationcategory{\mathcal{D}\!\mathit{ef}}
\usepackage{cleveref}
\crefname{originallemma}{original lemma}{original lemmas}
\Crefname{originallemma}{Original lemma}{Original lemmas}
\makeatletter
\newlabel{coherent-lemma-quasi-coherent-affine-cohomology-zero}{{01XB}{source}{}{sourcepointer-1}{}}
\newlabel{coherent-lemma-quasi-coherent-affine-cohomology-zero@cref}{{[originallemma][0][]01XB}{[1][0][]source}{}{}{}}
\makeatother
\crefname{originallemma}{original lemma}{original lemmas}
\Crefname{originallemma}{Original lemma}{Original lemmas}
\makeatletter
\newlabel{descent-proposition-same-cohomology-quasi-coherent}{{03DW}{source}{}{sourcepointer-2}{}}
\newlabel{descent-proposition-same-cohomology-quasi-coherent@cref}{{[originallemma][0][]03DW}{[1][0][]source}{}{}{}}
\makeatother
\crefname{originallemma}{original lemma}{original lemmas}
\Crefname{originallemma}{Original lemma}{Original lemmas}
\makeatletter
\newlabel{equation-fundamental-equivalence}{{05ZM}{source}{}{sourcepointer-3}{}}
\newlabel{equation-fundamental-equivalence@cref}{{[originallemma][0][]05ZM}{[1][0][]source}{}{}{}}
\makeatother
\crefname{originallemma}{original lemma}{original lemmas}
\Crefname{originallemma}{Original lemma}{Original lemmas}
\makeatletter
\newlabel{spaces-limits-lemma-affine}{{07VS}{source}{}{sourcepointer-4}{}}
\newlabel{spaces-limits-lemma-affine@cref}{{[originallemma][0][]07VS}{[1][0][]source}{}{}{}}
\makeatother
\crefname{originallemma}{original lemma}{original lemmas}
\Crefname{originallemma}{Original lemma}{Original lemmas}
\makeatletter
\newlabel{spaces-morphisms-lemma-i-star-equivalence}{{04CJ}{source}{}{sourcepointer-5}{}}
\newlabel{spaces-morphisms-lemma-i-star-equivalence@cref}{{[originallemma][0][]04CJ}{[1][0][]source}{}{}{}}
\makeatother
\crefname{originallemma}{original lemma}{original lemmas}
\Crefname{originallemma}{Original lemma}{Original lemmas}
\makeatletter
\newlabel{theorem-topological-invariance}{{05ZH}{source}{}{sourcepointer-6}{}}
\newlabel{theorem-topological-invariance@cref}{{[originallemma][0][]05ZH}{[1][0][]source}{}{}{}}
\makeatother
\title{Sections of the structure sheaf of a thickening surject onto those of the closed subspace on affine etale objects}
\author{The Stacks Project Authors}
\date{stacks-project:spaces-more-morphisms.tex:local-snapshot-20261011; GNU FDL 1.2 or later}
\begin{document}
\maketitle
\noindent\textit{Editorial note.} All delivered mathematical text is the author's own, copied from the pinned Stacks Project chapter file spaces-more-morphisms.tex; no statement or proof step was rewritten, reordered or removed.\par\medskip
\noindent\textit{Editorial note.} Preamble provenance, as required for this source: the first line of the pinned chapter file is \textbackslash{}input\{preamble\}, and the native preamble was not staged with the source, so no preamble text lives inside the chapter. The retained preamble used here is the project's single Stacks-wide preamble (the same 217-line / 8116-byte file that the already published chow and etale-cohomology cards use), copied mechanically line for line from sources/stacks-chow/preamble.tex lines 14-18 (Xy-pic setup) and 154-217 (theorem declarations and mathematical macros). Not one character was changed, no package was added, no macro was defined, and the excerpt serves only as preamble, never as a span coordinate.\par\medskip
\noindent\textit{Editorial note.} The same-chapter pointers are located at their \textbackslash{}label lines inside the pinned chapter file; the cross-chapter pointers are located in the project's own tag list staged in the snapshot (sources/stacks-chow/tags.txt), at the verbatim lines carrying their canonical tags.\par\medskip
\noindent\textit{Editorial note.} The author supplies two proof environments for this lemma, the second titled 'Proof for finite order thickenings' and introduced in the text as the case most used in practice (proof lines 1696-1706 and 1708-1727). Both are counted, in the order the author gives them, which is why this card carries two proof spans.\par\medskip
\noindent\textit{Allowed standard prerequisites (editorial).}\begin{itemize}
\item Etale localization of thickenings of algebraic spaces. Same chapter, Theorem labelled theorem-topological-invariance; invoked at line 1698. Produces the etale U' -> X' with U = X x\_\{X'\} U' for the affine object U of X\_\{spaces,etale\}.
\item Affineness descends along the thickening. Limits of Spaces, Lemma labelled spaces-limits-lemma-affine; invoked at lines 1699-1700. Shows that U' is an affine scheme, so its global sections can be compared with those of U.
\item The kernel of a first order thickening is a quasi-coherent module on the closed subspace. Morphisms of Spaces, Lemma labelled spaces-morphisms-lemma-i-star-equivalence; invoked at lines 1713-1716. Identifies the square-zero ideal I with a quasi-coherent O\_X-module in the second proof.
\item Etale cohomology of a quasi-coherent module on an affine scheme vanishes in positive degrees. Descent, Proposition labelled descent-proposition-same-cohomology-quasi-coherent and Cohomology of Schemes, Lemma labelled coherent-lemma-quasi-coherent-affine-cohomology-zero; invoked at lines 1722-1726. Kills H\textasciicircum{}1\_etale(U, I), which is what makes the section map surjective.
\end{itemize}\medskip
\ifdefined\hypertarget\hypertarget{sourcepointer-1}{}\fi
\noindent\textit{Source reference.} Original lemma 01XB: sources/stacks-chow/tags.txt, source lines 2512--2512. Navigation pointer for \textbackslash{}ref\{coherent-lemma-quasi-coherent-affine-cohomology-zero\} used inside the counted span (line 1726); the argument that consumes it is retained in the span.\par\medskip
\ifdefined\hypertarget\hypertarget{sourcepointer-2}{}\fi
\noindent\textit{Source reference.} Original lemma 03DW: sources/stacks-chow/tags.txt, source lines 4435--4435. Navigation pointer for \textbackslash{}ref\{descent-proposition-same-cohomology-quasi-coherent\} used inside the counted span (line 1724); the argument that consumes it is retained in the span.\par\medskip
\ifdefined\hypertarget\hypertarget{sourcepointer-3}{}\fi
\noindent\textit{Source reference.} Original lemma 05ZM: sources/stacks-spaces-more-morphisms/pinned-original-source.tex, source lines 1554--1554. Navigation pointer for \textbackslash{}ref\{equation-fundamental-equivalence\} used inside the counted span (line 1693); the argument that consumes it is retained in the span.\par\medskip
\ifdefined\hypertarget\hypertarget{sourcepointer-4}{}\fi
\noindent\textit{Source reference.} Original lemma 07VS: sources/stacks-chow/tags.txt, source lines 10017--10017. Navigation pointer for \textbackslash{}ref\{spaces-limits-lemma-affine\} used inside the counted span (line 1699); the argument that consumes it is retained in the span.\par\medskip
\ifdefined\hypertarget\hypertarget{sourcepointer-5}{}\fi
\noindent\textit{Source reference.} Original lemma 04CJ: sources/stacks-chow/tags.txt, source lines 5693--5693. Navigation pointer for \textbackslash{}ref\{spaces-morphisms-lemma-i-star-equivalence\} used inside the counted span (line 1715); the argument that consumes it is retained in the span.\par\medskip
\ifdefined\hypertarget\hypertarget{sourcepointer-6}{}\fi
\noindent\textit{Source reference.} Original lemma 05ZH: sources/stacks-spaces-more-morphisms/pinned-original-source.tex, source lines 1316--1316. Navigation pointer for \textbackslash{}ref\{theorem-topological-invariance\} used inside the counted span (line 1698); the argument that consumes it is retained in the span.\par\medskip
\begin{lemma}
\label{lemma-first-order-thickening-surjective}
Let $S$ be a scheme. Let $X \subset X'$ be a thickening
of algebraic spaces over $S$. Let $U$ be an affine object of
$X_{spaces, \etale}$. Then
$$
\Gamma(U, \mathcal{O}_{X'}) \to \Gamma(U, \mathcal{O}_X)
$$
is surjective where we think of $\mathcal{O}_{X'}$ as a sheaf on
$X_{spaces, \etale}$ via (\ref{equation-fundamental-equivalence}).
\end{lemma}

\begin{proof}
Let $U' \to X'$ be the \'etale morphism of algebraic spaces such that
$U = X \times_{X'} U'$, see Theorem \ref{theorem-topological-invariance}.
By Limits of Spaces, Lemma \ref{spaces-limits-lemma-affine} we see
that $U'$ is an affine scheme. Hence
$\Gamma(U, \mathcal{O}_{X'}) = \Gamma(U', \mathcal{O}_{U'}) \to
\Gamma(U, \mathcal{O}_U)$
is surjective as $U \to U'$ is a closed immersion of affine schemes.
Below we give a direct proof for finite order thickenings
which is the case most used in practice.
\end{proof}

\begin{proof}[Proof for finite order thickenings]
We may assume that $X \subset X'$ is a first order thickening by the
principle explained above. Denote $\mathcal{I}$ the kernel of the surjection
$\mathcal{O}_{X'} \to \mathcal{O}_X$. As $\mathcal{I}$ is a quasi-coherent
$\mathcal{O}_{X'}$-module and since $\mathcal{I}^2 = 0$ by the definition
of a first order thickening we may apply
Morphisms of Spaces, Lemma
\ref{spaces-morphisms-lemma-i-star-equivalence}
to see that $\mathcal{I}$ is a quasi-coherent $\mathcal{O}_X$-module.
Hence the lemma follows from the long exact cohomology sequence
associated to the short exact sequence
$$
0 \to \mathcal{I} \to \mathcal{O}_{X'} \to \mathcal{O}_X \to 0
$$
and the fact that $H^1_\etale(U, \mathcal{I}) = 0$ as
$\mathcal{I}$ is quasi-coherent, see
Descent, Proposition \ref{descent-proposition-same-cohomology-quasi-coherent}
and Cohomology of Schemes, Lemma
\ref{coherent-lemma-quasi-coherent-affine-cohomology-zero}.
\end{proof}

\end{document}
